\documentclass[10pt,a4paper]{amsart}
\usepackage[margin=19mm]{geometry}
\usepackage{amsmath,amssymb,mathtools}
\usepackage[scr=rsfs,cal=euler]{mathalfa}
\usepackage{xfrac}
\usepackage[unicode,hidelinks,bookmarks=false]{hyperref}
\newcommand{\ie}{i.e.\ }
\newcommand{\cf}{cf.\ }
\newcommand{\resp}{respectively\ }
\newcommand{\Subsection}{\subsection}
\providecommand{\nobreakdash}{\nobreak\hbox{-}}
\setlength{\emergencystretch}{2em}
\multlinegap0pt
\usepackage{bm}
\usepackage{bbm}
\usepackage{dsfont}
\usepackage{xspace}
\usepackage{cancel}
\usepackage{mathtools}
\usepackage{amsthm}
\makeatletter
\@ifundefined{theorem}{\newtheorem{theorem}{Theorem}}{}
\makeatother
\title[Multiscale analysis via pseudo-reversing and applications to]{Multiscale analysis via pseudo-reversing and applications to manifold-valued sequences}
\author{Wael Mattar, Nir Sharon}
\date{Source: arXiv:2305.06261v4 (CC BY 4.0)}
\begin{document}
\maketitle

\noindent\emph{Source and licence.} Multiscale analysis via pseudo-reversing and applications to manifold-valued sequences. Version arXiv:2305.06261v4 (CC BY 4.0), distributed under the Creative Commons Attribution 4.0 International licence, as recorded in the arXiv licence metadata for this exact version and shown on the abstract page https://arxiv.org/abs/2305.06261v4. Full rights notice: licenses/arxiv-2305.06261-CC-BY-4.0.txt.\par\smallskip

\noindent\textit{Editorial scope.} Counted unit: the Theorem whose source label is \texttt{thm:synthesis\_theorem} in the article (source0.tex lines 493--499, source label \texttt{thm:synthesis\_theorem}), delivered together with the article's own complete original proof of it (source0.tex lines 501--528, one \texttt{proof} environment of 28 lines). No mathematics is written, repaired or re-derived: statement and proof are the article's own text line for line. Required context retained uncounted: linear\_multiscale\_analysis (lines 252--254); linear\_multiscale\_synthesis (lines 256--258); linear\_subdivision\_scheme (lines 294--296); eqn:subdivision\_norm (lines 414--416); linear\_details\_bound (lines 431--433); differentiable\_details\_bound (lines 466--468). Every definition, display and label of those lines is retained unchanged. Omitted from this package and not redistributed: 1--251, 255--255, 259--293, 297--413, 417--430, 434--465, 469--492, 500--500, 529--1062. Nothing inside a retained excerpt is omitted or altered: every retained body is delivered exactly as the article's lines are written, so no omission receipt of any kind accompanies this package. Citations: the retained excerpts carry no citation command at all, so no bibliography entry of the article is transcribed and no reference is redistributed. Reference adaptation: none was needed and none was made; every reference inside the delivered unit resolves inside it, so no unresolved reference or citation is produced. Printed slips kept literal: no printed word, symbol, hypothesis or number of the counted statement or of its proof was altered. Layout and class: the article's own class is \texttt{article} with packages including amsfonts, amsmath, amssymb, amsthm, xcolor, color, floatrow, makeidx, graphicx, float, mathtools, geometry, subcaption, tikz-cd; the wrapper is the lane's pinned \texttt{amsart} editorial wrapper at 10pt on A4, which restates the theorem-like environments the retained text needs and adds the article's own macro definitions that the retained text needs (none), with extra wrapper packages (bm, bbm, dsfont, xspace, cancel, mathtools). The delivered numbering is the unit's own. Assets: no retained span contains a figure environment, a graphics-inclusion command, a PDF-page inclusion, a table or a TikZ picture; no image file is copied into this package, the package declares no asset dependency and no image file is redistributed with it. Licence: version arXiv:2305.06261v4 of this article is distributed under the Creative Commons Attribution 4.0 International licence, as recorded in the arXiv licence metadata for this exact version and shown on the abstract page https://arxiv.org/abs/2305.06261v4. A full rights notice is delivered at licenses/arxiv-2305.06261-CC-BY-4.0.txt. Characters: every retained excerpt is pure ASCII, so the delivered engine, XeLaTeX with its default Latin Modern text font, reports no missing character.
\begin{equation}~\label{linear_multiscale_analysis}
\boldsymbol{c}^{(\ell-1)} = \mathcal{D}\boldsymbol{c}^{(\ell)}, \quad \boldsymbol{d}^{(\ell)} = \boldsymbol{c}^{(\ell)} - \mathcal{S}\boldsymbol{c}^{(\ell-1)}, \quad \ell=1,\dots,J,
\end{equation}

\begin{equation}~\label{linear_multiscale_synthesis}
\boldsymbol{c}^{(\ell)} = \mathcal{S}\boldsymbol{c}^{(\ell - 1)} + \boldsymbol{d}^{(\ell)}, \quad \ell=1,\dots,J.
\end{equation}

\begin{equation}~\label{linear_subdivision_scheme}
    \mathcal{S}_{\boldsymbol{\alpha}}(\boldsymbol{c})_j = \sum_{k\in\mathbb{Z}}\alpha_{j-2k}c_k, \quad j\in\mathbb{Z}.
\end{equation}

    \begin{equation}~\label{eqn:subdivision_norm}
        \|\mathcal{S}_{\boldsymbol{\beta}}\|_{\infty} = \max\big\{ \sum_{k\in\mathbb{Z}}|\beta_{2k}|, \; \sum_{k\in\mathbb{Z}}|\beta_{2k+1}|\big\}.
    \end{equation}

    \begin{equation}~\label{linear_details_bound}
        \|\boldsymbol{d}^{(\ell)}\downarrow 2\|_\infty \leq L(\xi)\Delta\boldsymbol{c}^{(\ell)},
    \end{equation}

    \begin{equation}~\label{differentiable_details_bound}
        \|\boldsymbol{d}^{(\ell)}\downarrow 2\|_\infty \leq P(\xi)(2\|\boldsymbol{\gamma}\|_1)^{-\ell},
    \end{equation}

\begin{theorem}~\label{thm:synthesis_theorem}
    Let $\boldsymbol{c}^{(J)}$ be a real-valued sequence, and let $\mathcal{S}_{\boldsymbol{\alpha}}$ be a refinement operator~\eqref{linear_subdivision_scheme} with a pseudo-reverse $\mathcal{D}_{\boldsymbol{\gamma}}$ for some $\xi>0$. Denote by $\{\boldsymbol{c}^{(0)};\boldsymbol{d}^{(1)},\dots,\boldsymbol{d}^{(J)}\}$ the multiscale representation~\eqref{linear_multiscale_analysis} of $\boldsymbol{c}^{(J)}$ using $\mathcal{S}_{\boldsymbol{\alpha}}$ and $ \mathcal{D}_{\boldsymbol{\gamma}}$. Then, there exists $C>0$ such that
    \begin{equation}~\label{synthesis_bound}
        \|\boldsymbol{c}^{(J)} - \boldsymbol{\zeta}^{(J)}\|_\infty \leq C\sum_{\ell=1}^{J}\|\boldsymbol{d}^{(\ell)}\downarrow 2\|_\infty,
    \end{equation}
    where $\boldsymbol{\zeta}^{(J)}$ is the synthesized sequence after setting the even detail coefficients of the multiscale representation to zero. In other words, the reconstruction error is proportional to the cumulative errors studied in~\eqref{linear_details_bound} and~\eqref{differentiable_details_bound}.
\end{theorem}

\begin{proof}
    We first note that $\boldsymbol{c}^{(J)}$ can be fully recovered via~\eqref{linear_multiscale_synthesis} when considering \emph{all} the detail coefficients of its pyramid representation. In other words, the iterations~\eqref{linear_multiscale_analysis} and~\eqref{linear_multiscale_synthesis} are algebraically compatible. In particular,
    \begin{equation*}
        \boldsymbol{c}^{(J)}=\mathcal{S}_{\boldsymbol{\alpha}}^{J}\boldsymbol{c}^{(0)} + \sum_{\ell=1}^J\mathcal{S}_{\boldsymbol{\alpha}}^{J-\ell}\boldsymbol{d}^{(\ell)},
    \end{equation*}
    where $\mathcal{S}_{\boldsymbol{\alpha}}^0=\mathcal{I}$ is the identity operator.

    Now, we consider the pyramid $\{\boldsymbol{c}^{(0)};\boldsymbol{q}^{(1)},\dots,\boldsymbol{q}^{(J)}\}$ where
    \begin{equation*}
        q^{(\ell)}_{j}=\begin{cases}
            d^{(\ell)}_j, & j=2n+1, \\
            0, & j=2n,
            \end{cases} \quad \quad\ell=1,\dots, J, \quad j\in\mathbb{Z},
    \end{equation*}
    and study its synthesis via~\eqref{linear_multiscale_synthesis} using $\mathcal{S}_{\boldsymbol{\alpha}}$. Our objective now is to quantify the distance (induced by the $\ell_\infty$ norm) between the reconstruction and the original input $\boldsymbol{c}^{(J)}$.

    To this end, we denote by 
    \begin{equation*}
        \boldsymbol{\zeta}^{(J)}=\mathcal{S}_{\boldsymbol{\alpha}}^{J}\boldsymbol{c}^{(0)} + \sum_{\ell=1}^J\mathcal{S}_{\boldsymbol{\alpha}}^{J-\ell}\boldsymbol{q}^{(\ell)},
    \end{equation*}
    and observe that
    \begin{equation*}
        \|\boldsymbol{c}^{(J)}-\boldsymbol{\zeta}^{(J)}\|_\infty \leq \sum_{\ell=1}^J\|\mathcal{S}_{\boldsymbol{\alpha}}^{J-\ell}(\boldsymbol{d}^{(\ell)}-\boldsymbol{q}^{(\ell)})\|_\infty \leq \sum_{\ell=1}^J\|\mathcal{S}_{\boldsymbol{\alpha}}^{J-\ell}\|_{\infty}\|\boldsymbol{d}^{(\ell)}-\boldsymbol{q}^{(\ell)}\|_\infty,
    \end{equation*}
    where $\|\mathcal{S}_{\boldsymbol{\alpha}}^{J-\ell}\|_{\infty}$ is the operator norm of the refinement $\mathcal{S}_{\boldsymbol{\alpha}}^{J-\ell}$. See~\eqref{eqn:subdivision_norm} for an explicit formula of the norm.

    The uniform boundedness principle guarantees the existence of a constant $C>0$ such that $\sup_{j\in\mathbb{N}} \|\mathcal{S}_{\boldsymbol{\alpha}}^j\|_{\infty}\leq C$. Now, notice that $\boldsymbol{d}^{(\ell)}$ and $\boldsymbol{q}^{(\ell)}$ may differ only on even indices, and hence $\|\boldsymbol{d}^{(\ell)}-\boldsymbol{q}^{(\ell)}\|_\infty=\|\boldsymbol{d}^{(\ell)}\downarrow 2\|_\infty$. We then immediately get the desired result~\eqref{synthesis_bound}.
\end{proof}

\end{document}
